\section{Conclusion}
\label{sec:conclusion}

On MQAR at length 2048 with a fixed fine budget, a two-level block selector recalled more as it kept more
super-blocks, up to a point: on identical data, keeping six beat keeping two on all four seeds, and keeping all twelve
(effectively flat selection) did worse on the one seed where both were run. Where dense attention learned, it beat
both settings run on that seed, by 0.291 at the cheapest and 0.055 at six super-blocks, so none met our pre-registered bar of 0.02.
Because this task lets the coarse level be solved by position, the curve speaks to how reliably the selector learns to
choose, not to information lost by pruning. Replaying the trained models, errors followed wrong block choices rather than failed
reading; the cheapest setting kept missing the right super-block despite the shortcut, and two intuitive repairs made
recall worse. Dense attention itself learned on only two of four seeds
within our training budget. These results come from one synthetic task with a positional shortcut, small models and few
seeds. We hope the protocol---matched data, reserved seeds, and every failure reported---is useful to others
measuring what sparse attention gives up.
